\section{Prior shift}
\label{sec:prior-shift}

\subsection{Correcting for a new spam rate (T5)}
\label{sec:t5}
% Member 5

A filter is trained where a share $\pio$ of emails is spam, but at deployment the share may be $\pion$.
Suppose spam still \emph{looks} the same: $p(x\mid y)$ is unchanged and only the class proportions move.
This is called label shift. By Bayes' rule the posterior odds are the likelihood ratio times the prior odds,
\begin{equation}
  \frac{\Prob(y=1\mid x)}{\Prob(y=0\mid x)} = \frac{p(x\mid y=1)}{p(x\mid y=0)}\cdot\frac{\pio}{\piz},
  \label{eq:t5-odds-train}
\end{equation}
and at deployment the same holds with $\pion,\pizn$ and the \emph{same} likelihood ratio. Dividing the two
removes the unknown likelihood ratio:
\begin{equation}
  \boxed{\;\frac{\Prob'(y=1\mid x)}{\Prob'(y=0\mid x)}
  = \frac{\Prob(y=1\mid x)}{\Prob(y=0\mid x)}\cdot\frac{\pion/\pizn}{\pio/\piz}\;}
  \label{eq:t5-odds}
\end{equation}
So we multiply the model's odds by the new prior odds over the old ones; $p(x\mid y)$ is never needed.
Here $\post=\Prob(y=1\mid x)$ is the model's output (written $q$ in Section~\ref{sec:calibration}).
Taking logarithms turns the product into a sum. With $a(x)=\logit\post(x)=\log\frac{\post(x)}{1-\post(x)}$,
the log posterior odds of the lectures,
\begin{equation}
  a'(x) = a(x) + \lshift, \qquad
  \lshift = \log\frac{\pion}{\pizn} - \log\frac{\pio}{\piz}, \qquad
  \post'(x) = \sigm\big(\logit\post(x) + \lshift\big),
  \label{eq:t5-shift}
\end{equation}
so the log-odds of every email move by the same amount $\lshift$.

\emph{Example.} Trained at $\pio=0.5$ and deployed at $\pion=0.1$, all odds are divided by $9$. An email
with $\post=0.98$ (odds $49$) gets $\post'\approx0.845$, so with the thresholds $\thz=0.3$, $\tho=0.97$ of
Section~\ref{sec:decision} it is inspected instead of sent to the spam folder.

\paragraph{In our two models.}
For logistic regression, $a(x)=w^\top\phi(x)$ with $\phi_0=1$, so adding $\lshift$ only replaces $w_0$ by
$w_0+\lshift$. Which $\pio$ did the model learn? At the fitted weights the gradient \eqref{eq:t3-grad} is zero. Its
component for $\phi_0=1$ is $\sum_n(\post_n-y_n)$, so $\sum_n\post_n=\sum_n y_n=N_1$: the average fitted
probability equals the training spam share $N_1/N$, which is the $\pio$ to use (exactly, since our intercept
is not penalised). For the Gaussian model the priors enter $a(x)$ only through
$\log(\pio/\piz)$ (Section~\ref{sec:t2}), for shared and for class-specific covariances, so replacing them by
$\pion,\pizn$ changes $a(x)$ by exactly $\lshift$: the correction is ``use the new priors''.

\paragraph{Keep the probability, move the thresholds.}
Since $\sigm$ is increasing,
\[
  \post'\ge\theta_k \iff \logit\post+\lshift\ge\logit\theta_k \iff \post\ge\sigm(\logit\theta_k-\lshift).
\]
With $r=e^{\lshift}$ the thresholds on the uncorrected $\post$ therefore become
\begin{equation}
  \theta_k^{\text{shift}} = \frac{\theta_k}{\theta_k + r\,(1-\theta_k)}, \qquad k=0,1,
  \label{eq:t5-thresholds}
\end{equation}
which gives the same action for every email (up to rounding exactly at a threshold). In the example
($r=1/9$) they move from $(0.30,\,0.97)$ to $(0.79,\,0.997)$; Figure~\ref{fig:t5} shows both views. Moving
thresholds leaves the model untouched; correcting $\post$ is needed when the probability itself is used.

\begin{figure}[ht]
\centering
\includegraphics[width=\linewidth]{t5_prior_correction.pdf}
\caption{Prior-shift correction for a model trained at $\pio=0.5$. Left: corrected $\post'$ against $\post$
\eqref{eq:t5-shift}. Right: the equivalent thresholds on the uncorrected $\post$ \eqref{eq:t5-thresholds}
for every deployment share $\pion$; the shaded band is the inspect region.}
\label{fig:t5}
\end{figure}

If the trained posterior is exact, \eqref{eq:t5-odds} gives the exact posterior at the new prior, so the
corrected decisions are the Bayes decisions for the deployment population and no rule does better there.
Section~\ref{sec:c5} measures how much the uncorrected filter loses.

\begin{assumptions}
\begin{enumerate}[label=(\roman*),itemsep=1pt,topsep=2pt]
  \item \emph{Label shift only:} if spam changes its wording, the likelihood ratio changes and no single
        $\lshift$ can fix it.
  \item \emph{Calibrated training probabilities:} the correction multiplies the model's odds, so any error
        in them carries over. A constant shift cannot fix an over-confident model (that needs a change of
        slope, Section~\ref{sec:c4}), so calibrate first and correct for the prior second.
  \item \emph{Known priors:} $\pio$ is the training share, and $\pion$ must be known or estimated.
  \item Both priors lie strictly between $0$ and $1$.
\end{enumerate}
\end{assumptions}
\texttt{tests/test\_prior\_shift.py} checks that the correction reproduces the exact Gaussian posterior at
each new prior, that refitting logistic regression at a new prior moves only the intercept by $\lshift$, and
that the threshold form gives the same actions. The functions are in \texttt{prior\_shift.py}.

\subsection{Experiments (C5)}
\label{sec:c5}
% Member 1. Every number below comes from python -m project4.experiments.task4_prior_shift
% (results/c5_summary.csv, results/c5_em_priors.csv).

We now test the correction of Section~\ref{sec:t5} on data where only the class proportions change.
Throughout, $\cfp=10$, $\cfn=1$ and $\ci=0.3$.

\paragraph{Setup.}
\begin{itemize}[itemsep=1pt,topsep=2pt]
  \item \emph{Synthetic.} The two 5-dimensional Gaussian classes of Section~\ref{sec:task2-results}
        (different covariances). Training and validation data are drawn with $\pio=0.5$. For each
        deployment share $\pion$ in $\{0.05,\allowbreak 0.1,\allowbreak 0.2,\allowbreak 0.4,\allowbreak 0.5,\allowbreak
        0.6,\allowbreak 0.8\}$ we draw a fresh sample of $5{,}000$ emails at $\pion$, so $p(x\mid y)$ is exactly the
        same as in training. The share $0.5$ is the training share itself: a no-shift check.
  \item \emph{Spambase.} The usual split ($\pio=0.394$). Each deployment set is the test part resampled
        within each class to the share $\pion$, without replacement and as large as the test part allows
        (for example $586$ emails, $29$ of them spam, at $\pion=0.05$).
  \item \emph{Models.} The Gaussian classifier with shared covariance and logistic regression, each tuned on
        validation log-loss. Spambase features are transformed by $\log(1+x)$ before standardising, which gives
        the Gaussian model much better probabilities (Section~\ref{sec:models}).
  \item \emph{Decisions.} The Bayes rule of Section~\ref{sec:decision} applied to four versions of the posterior:
        \textbf{none} (the model's $\post$, trained at $\pio$);
        \textbf{known} ($\post'$ from \eqref{eq:t5-shift} with the true $\pion$);
        \textbf{EM} ($\post'$ with $\pion$ estimated from the unlabelled deployment emails, see below);
        and, on synthetic data only, the \textbf{oracle}, the true posterior at $\pion$.
\end{itemize}
Everything is repeated over the 20 splits; we report the mean realized loss with its standard error.

\paragraph{Estimating an unknown deployment prior.}
In practice $\pion$ is rarely known. Following Saerens et al.~\cite{saerens2002}, we estimate it from the
model's outputs on the $M$ unlabelled deployment emails by iterating
\[
  \hat{\pion} \;\leftarrow\; \frac1M\sum_{m=1}^M \post'_m(\hat{\pion}),
  \qquad \post'_m(\hat{\pion}) = \text{the corrected posterior \eqref{eq:t5-shift} of email $m$ at prior }\hat{\pion},
\]
starting from $\hat{\pion}=\pio$. This is the EM algorithm for the unknown class proportion, so its fixed point is
the maximum-likelihood estimate of $\pion$ given the model's posteriors. That estimate is close to the true
$\pion$ only if those posteriors are calibrated.

\begin{figure}[ht]
\centering
\includegraphics[width=\linewidth]{c5_prior_shift.pdf}
\caption{Realized loss per email against the deployment spam share $\pion$, for decisions on the uncorrected
posterior, on the posterior corrected with the true $\pion$, and on the posterior corrected with $\pion$
estimated by EM. Dashed: decisions on the true posterior (synthetic data only). Dotted line: the training share.
Bands show $\pm1$ standard error over 20 splits. Generated by \texttt{experiments/task4\_prior\_shift.py}.}
\label{fig:c5}
\end{figure}

\begin{table}[ht]
\centering
\small
\resizebox{\linewidth}{!}{\input{\resultsdir/c5_loss_synthetic.tex}}\\[6pt]
\resizebox{0.8\linewidth}{!}{\input{\resultsdir/c5_loss_spambase.tex}}
\caption{Mean realized loss ($\pm$ standard error over 20 splits) after a prior shift. Top: synthetic data
(training share $0.5$). Bottom: Spambase (training share $0.394$). Columns: no correction, correction with the
known $\pion$, correction with $\pion$ estimated by EM, and the true posterior.}
\label{tab:c5-loss}
\end{table}

\begin{table}[ht]
\centering
\small
\input{\resultsdir/c5_gain.tex}
\caption{Loss saved by the correction with the known $\pion$: uncorrected minus corrected loss, paired by
split, mean $\pm$ standard error over 20 splits. Pairing removes the variation the two methods share, so
these errors are much smaller than those of Table~\ref{tab:c5-loss}.}
\label{tab:c5-gain}
\end{table}

\begin{table}[ht]
\centering
\small
\resizebox{\linewidth}{!}{\input{\resultsdir/c5_confusion.tex}}
\caption{Mean confusion matrices per deployment set at $\pion=0.05$ (synthetic: $5{,}000$ emails;
Spambase: $586$ emails), without and with the correction (known $\pion$).}
\label{tab:c5-confusion}
\end{table}

\paragraph{Results.}
Figure~\ref{fig:c5} and Tables~\ref{tab:c5-loss} and~\ref{tab:c5-gain} show the following.
\begin{enumerate}[itemsep=2pt,topsep=2pt]
  \item \emph{The correction pays off whenever the share moves, and the gain grows with the shift.}
        At $\pion=0.05$ it cuts the loss by more than half for both models on both data sets: on Spambase from
        $0.095$ to $0.038$ (logistic) and from $0.108$ to $0.051$ (Gaussian), and on synthetic data from
        $0.154$ and $0.137$ to $0.040$, close to the oracle's $0.034$. At $\pion=0.8$ it cuts the loss by a
        quarter to a third. The paired gains of Table~\ref{tab:c5-gain} are positive by several standard errors
        whenever the share moves by $0.1$ or more on synthetic data, or by $0.2$ or more on Spambase; the small move
        from $0.394$ to $0.5$ on Spambase gains $0.003\pm0.002$ and $0.001\pm0.002$, within noise. Even on Spambase at
        $\pion=0.6$, where the bands of Figure~\ref{fig:c5} overlap,
        the gain is $0.012\pm0.003$ (logistic) and $0.015\pm0.003$ (Gaussian). At the training share the
        curves coincide, as they should: the paired gain is $-0.0001\pm0.0006$ and $0.0004\pm0.0002$ on Spambase at
        $\pion=0.4$ (training share $0.394$), and at most $0.0001$ on synthetic data at $\pion=0.5$.
  \item \emph{Where the gain comes from.} When spam becomes rare, the uncorrected filter still believes that
        $39\%$ (Spambase) or $50\%$ (synthetic) of emails are spam. It inspects far too many emails (on synthetic
        data $43$--$49\%$ of all emails, against $1.6\%$ after correction) and sends more real email to the
        spam folder (Table~\ref{tab:c5-confusion}: on Spambase $3.8$ against $0.8$ real emails per deployment set
        for logistic regression). The corrected filter instead delivers more spam, which costs only $\cfn=1$
        per email.
  \item \emph{Lower accuracy, lower cost.} The corrected decisions are slightly \emph{less} accurate at
        $\pion=0.05$ (synthetic: $0.964$ against $0.987$ on the emails not inspected) but three to four times
        cheaper. Much of the uncorrected filter's accuracy is bought by handing almost half of the emails to
        an inspector, at $\ci=0.3$ each. This is the effect of Section~\ref{sec:accuracy-cost}: the right target is the loss,
        not the error count.
  \item \emph{The remaining gap is the model, not the prior.} On synthetic data the corrected curves stay
        $0.04$--$0.06$ above the oracle near the training share (and $0.006$ above it at $\pion=0.05$, where
        almost every email is easy). The two classes have different covariances,
        so the true log-odds are quadratic in $x$, while both of our models have linear log-odds. Label shift
        holds exactly here by construction; what fails is the premise of Section~\ref{sec:t5} that the trained
        posterior is exact (assumption~(ii)). The prior correction cannot remove this error.
  \item \emph{An unknown prior can be estimated well if the model is calibrated.} On Spambase, logistic
        regression's EM estimate has a mean bias of at most $0.002$ and a mean absolute error of $0.007$--$0.013$
        at every $\pion$, and the EM-corrected loss is almost the same as with the known prior. The Gaussian model's
        estimate is too low: on synthetic data by $0.04$--$0.07$ at every share, even at the training share $0.5$
        (estimate $0.446$), and on Spambase by $0.036$--$0.047$ for $\pion\ge0.4$, already $0.364$ at
        $\pion=0.4$ where there is essentially no shift. A reliability check explains the direction. On
        $200{,}000$ fresh synthetic emails at the training share (split $0$), emails both models give
        $\post\approx0.06$ are $21$--$22\%$ spam, so they are over-confident at the low end, while emails given
        $\post\approx0.85$ are $93$--$94\%$ spam, under-confident at the high end
        (\texttt{results/c5\_calibration\_synthetic.csv}). Averaging such probabilities gives too small a share,
        most of all when most emails sit at the low end: logistic regression estimates $0.015$ instead of $0.05$.
        The loss rises only a little ($0.046$ against $0.040$). This is assumption~(ii) of Section~\ref{sec:t5}
        at work: EM inherits any miscalibration of the model.
\end{enumerate}

\paragraph{Caveats.}
The Spambase deployment sets are resampled from a test part of $919$ emails, so at extreme shares they are
small and reuse the same emails, and their error bands are wider. All 20 Spambase splits are drawn from
the same $4{,}601$ emails, so they are not independent and their standard errors are somewhat optimistic.
Resampling within each class keeps
$p(x\mid y)$ fixed by construction. A real change of spam content would also change $p(x\mid y)$, and no
prior correction can fix that.

\begin{assumptions}
Label shift only (enforced by the sampling); the training share $\pio$ is the share in the training part;
the correction uses the model's posteriors as if they were calibrated; EM uses only unlabelled deployment emails.
\end{assumptions}
